%% Appendix B — MQTT Topic Reference
\chapter{MQTT Topic Reference}
\label{app:mqtt_reference}

\section{Topic Naming Convention}

All MMS V2 topics use the prefix \texttt{mms/}. The hierarchy is:
\begin{verbatim}
mms/
  nodes/
    {machine_id}/
      status      — node publishes live state
      event       — node publishes completed session events
      heartbeat   — node publishes health data
      command     — bridge publishes commands to node
      config      — bridge publishes config to node (retained)
  all/
    revoked       — bridge publishes revocation list (retained)
\end{verbatim}

\section{Topic Details}

\subsection{\texttt{mms/nodes/\{id\}/status}}
\begin{itemize}[leftmargin=1.5em,noitemsep]
    \item \textbf{Publisher}: ESP32 access node
    \item \textbf{QoS}: 1
    \item \textbf{Retained}: Yes
    \item \textbf{Last Will}: \texttt{\{"state":"offline"\}} (retained; broker publishes on unexpected disconnect)
\end{itemize}

\begin{lstlisting}[style=jsonlstyle,caption={Status topic payload}]
{
  "state":            "active",
  "roll_number":      "220002123",
  "session_start":    1718600000,
  "firmware_version": "2.0.0",
  "ip_address":       "192.168.0.101",
  "last_seen":        1718600060,
  "rssi":             -55
}
\end{lstlisting}

Valid \texttt{state} values: \texttt{"offline"}, \texttt{"idle"}, \texttt{"card\_present"}, \texttt{"active"}, \texttt{"error"}, \texttt{"ota"}.

\subsection{\texttt{mms/nodes/\{id\}/event}}
\begin{itemize}[leftmargin=1.5em,noitemsep]
    \item \textbf{Publisher}: ESP32 access node
    \item \textbf{QoS}: 1
    \item \textbf{Retained}: No
\end{itemize}

\begin{lstlisting}[style=jsonlstyle,caption={Event topic payload}]
{
  "t":  1718600900,
  "r":  "220002123",
  "m":  3,
  "e":  "session_end",
  "d":  245,
  "er": "card_removed"
}
\end{lstlisting}

Fields: \texttt{t}=end timestamp, \texttt{r}=roll number, \texttt{m}=machine\_id, \texttt{e}=event type, \texttt{d}=duration\_s, \texttt{er}=end\_reason.

\subsection{\texttt{mms/nodes/\{id\}/heartbeat}}
\begin{itemize}[leftmargin=1.5em,noitemsep]
    \item \textbf{Publisher}: ESP32 access node (every 10 seconds)
    \item \textbf{QoS}: 0
    \item \textbf{Retained}: No
\end{itemize}

\begin{lstlisting}[style=jsonlstyle,caption={Heartbeat topic payload}]
{"ts": 1718600060, "heap_free": 218432, "uptime_s": 3600}
\end{lstlisting}

\subsection{\texttt{mms/nodes/\{id\}/command}}
\begin{itemize}[leftmargin=1.5em,noitemsep]
    \item \textbf{Publisher}: RPi bridge (forwards from RTDB)
    \item \textbf{QoS}: 1
    \item \textbf{Retained}: No
\end{itemize}

\begin{lstlisting}[style=jsonlstyle,caption={Command topic payload}]
{
  "cmd":       "stop",
  "issued_by": "admin_uid",
  "ts":        1718600100,
  "url":       null,
  "version":   null
}
\end{lstlisting}

Valid \texttt{cmd} values: \texttt{"stop"}, \texttt{"reboot"}, \texttt{"ota"}.

For OTA commands: \texttt{"url"} contains the firmware download URL, \texttt{"version"} contains the target version string.

\subsection{\texttt{mms/nodes/\{id\}/config}}
\begin{itemize}[leftmargin=1.5em,noitemsep]
    \item \textbf{Publisher}: RPi bridge (from Firestore \texttt{onSnapshot})
    \item \textbf{QoS}: 1
    \item \textbf{Retained}: Yes (nodes receive latest config on connect even if bridge is not running)
\end{itemize}

\begin{lstlisting}[style=jsonlstyle,caption={Config topic payload}]
{
  "name":               "Soldering Station 1",
  "session_limit_min":  15,
  "relay_active_high":  true,
  "machine_active":     true
}
\end{lstlisting}

\subsection{\texttt{mms/all/revoked}}
\begin{itemize}[leftmargin=1.5em,noitemsep]
    \item \textbf{Publisher}: RPi bridge (from Firestore \texttt{onSnapshot})
    \item \textbf{QoS}: 1
    \item \textbf{Retained}: Yes (all nodes receive on connect; no need to re-push)
\end{itemize}

\begin{lstlisting}[style=jsonlstyle,caption={Revocation list topic payload}]
{
  "uids":       ["A1B2C3D4", "E5F60718"],
  "updated_at": 1718600300
}
\end{lstlisting}

\section{Subscribing to Topics for Debugging}

\begin{lstlisting}[style=bashstyle]
# On the RPi or any machine with mosquitto-clients installed

# Watch all MMS traffic
mosquitto_sub -h 192.168.0.10 -t "mms/#" -v

# Watch only status from all nodes
mosquitto_sub -h 192.168.0.10 -t "mms/nodes/+/status" -v

# Watch events from node 3
mosquitto_sub -h 192.168.0.10 -t "mms/nodes/3/event" -v

# Watch revocation list
mosquitto_sub -h 192.168.0.10 -t "mms/all/revoked"

# Manually publish a stop command (for testing)
mosquitto_pub -h 192.168.0.10 -t "mms/nodes/3/command" \
  -m '{"cmd":"stop","issued_by":"test","ts":1718600100}'
\end{lstlisting}
